Figure 1 is a whole-sky star chart of Yanqing, Beijing at 20:30 tonight
(UTC+8) with the limiting magnitude $= 5^{\mathrm{m}}$ ($\mathrm{m}$ =
magnitude). Four stars (about $1^{\mathrm{m}}$--$3^{\mathrm{m}}$) and one
planet (brighter than $2^{\mathrm{m}}$) are missing in this chart. In the
chart, the distance from the centre is in proportion to zenith distance.

% FIGURE NEEDED: Figure 1, whole-sky star chart of Yanqing, Beijing at 20:30
% (2018_DAY.pdf page 3)

\begin{description}
\item[(1)] (20 points) Draw a cross (X) on the location of each missing star
  and mark ``T'' on the chart, and draw a cross (X) on the location of the
  missing planet and mark ``P'' on the chart.
\item[(2)] (5 points) Please mark the orientation of the star chart with ``N''
  ``E'' ``S'' ``W'' at the edge of the star chart.
\item[(3)] (10 points) On the chart, the celestial equator passes through many
  constellations. Please write down the name of any five of these
  constellations (IAU codes).
\item[(4)] (5 points) Using the star chart, estimate the altitude of Aldebaran
  ($\alpha$ Tau), to the nearest degree.
\end{description}
